\documentclass[11pt]{article}
\usepackage{amsmath,amssymb,amsthm}
\usepackage[margin=1in]{geometry}
\usepackage{hyperref}
\title{Disproof of TLMC Conjecture 00000001168}
\author{}
\date{October 2, 2026}
\newtheorem{theorem}{Theorem}
\newtheorem{lemma}[theorem]{Lemma}
\newcommand{\SSp}{\mathbb{S}}
\newcommand{\R}{\mathbb{R}}
\begin{document}
\maketitle

\section*{Verdict}
\textbf{FALSE.} The round metric on $\mathrm{SU}(2) \cong \SSp^3$ is itself a
Zoll metric, and its first nonzero Laplace eigenvalue has multiplicity
$4 > 2$; the family of rotated round metrics that the conjecture itself
cites violates the claimed bound on every member.

\section{The conjecture}

\begin{quote}
\textit{Definition: A Zoll metric is one all of whose geodesics are closed.
Conjecture: The spectral multiplicity upper bound for Zoll metrics on
$\mathrm{SU}(2)$ is two; and the counterexample to multiplicity rigidity is
the family of rotated round metrics. (Zoll multiplicity rigidity)}
\end{quote}

\section{Counterexample: the round metric on \texorpdfstring{$\mathrm{SU}(2)$}{SU(2)}}

Identify $\mathrm{SU}(2)$ with the unit quaternions, i.e.\ with the round
unit sphere $\SSp^3 \subset \R^4$. Write $g_r$ for the round (constant
curvature $1$) metric.

\begin{lemma}\label{lem:zoll}
$(\SSp^3, g_r)$ is a Zoll metric in the sense of the conjecture's own
definition: every geodesic is a great circle
$\gamma(t) = \cos(t)\, u + \sin(t)\, v$ with $u, v \in \SSp^3$ orthonormal,
closed with the common period $2\pi$.
\end{lemma}

\begin{proof}
The geodesics of the round sphere are exactly the great circles
$C = \SSp^3 \cap P$, $P \subset \R^4$ a two-plane through the origin, and the
parametrization above (unit speed) satisfies the geodesic equation
$\ddot\gamma + \gamma = 0$ identically; since $\cos$ and $\sin$ have common
period $2\pi$, $\gamma(t + 2\pi) = \gamma(t)$ for every great circle. The
script \texttt{reproduce.py} verifies $|\gamma(t)| = 1$,
$\ddot\gamma + \gamma = 0$ at sample points, and the $2\pi$-periodicity for
$200$ sampled orthonormal pairs $(u,v)$.
\end{proof}

\begin{theorem}\label{thm:main}
The $k$-th eigenvalue of the Laplace--Beltrami operator of $(\SSp^3, g_r)$ is
$\lambda_k = k(k+2)$, $k = 0, 1, 2, \dots$, with multiplicity $(k+1)^2$. In
particular $\lambda_1 = 3$ has multiplicity $4 > 2$. Hence the round metric
$g_r$, transported to $\mathrm{SU}(2)$ along the quaternion identification,
is a Zoll metric on $\mathrm{SU}(2)$ whose spectrum violates the claimed
upper bound of two. The conjecture is false.
\end{theorem}

\begin{proof}
The eigenfunctions on the round $\SSp^3 \subset \R^4$ are the restrictions of
the homogeneous harmonic polynomials on $\R^4$: the eigenspace for
$\lambda_k$ is the restriction of $\mathcal{H}_k$, the kernel of the flat
Laplacian $\Delta$ restricted to the space $\mathrm{Hom}_k$ of homogeneous
polynomials of degree $k$ in $4$ variables.

\emph{Dimension count.} $\dim \mathrm{Hom}_k$ is the number of monomials of
degree $k$ in $4$ variables,
\[
\dim \mathrm{Hom}_k = \binom{k+3}{3} = \frac{(k+1)(k+2)(k+3)}{6}.
\]
Since $\Delta$ maps $\mathrm{Hom}_k$ onto $\mathrm{Hom}_{k-2}$ (classical;
$\Delta$ is surjective because it is self-adjoint for the Fischer inner
product and its image is the orthogonal complement of the harmonic space in
$\mathrm{Hom}_{k-2}$; surjectivity is also verified computationally below),
and $\dim \mathrm{Hom}_{k-2} = \binom{k+1}{3}$, we get
\[
\dim \mathcal{H}_k = \binom{k+3}{3} - \binom{k+1}{3} = (k+1)^2 .
\]

\emph{The case $k=1$ explicitly.} $\mathrm{Hom}_1$ is the space of linear
forms $c_1 x_1 + c_2 x_2 + c_3 x_3 + c_4 x_4$, all of which are harmonic;
so $\lambda_1 = 1 \cdot 3 = 3$ and the eigenspace contains the restrictions
of the four coordinate functions. These are linearly independent: the
evaluation matrix at the four axis points $e_1, \dots, e_4 \in \SSp^3$ is the
identity ($x_j(e_i) = \delta_{ij}$), so a linear combination vanishing on
$\SSp^3$ has $c_i = 0$ for every $i$. Hence
$\dim \mathcal{H}_1 = 4$, matching $(1+1)^2 = 4$.

\emph{Independent computational verification.} The script
\texttt{reproduce.py} assembles the flat Laplacian
$\Delta\colon \mathrm{Hom}_k \to \mathrm{Hom}_{k-2}$ as an explicit matrix
over $\mathbb{Q}$ (entry
$\partial^2/\partial x_i^2\, x^a = a_i(a_i - 1)\, x^{a - 2e_i}$) and computes
its rank by exact rational Gaussian elimination for $k = 0, \dots, 8$. In
every case the map is surjective onto $\mathrm{Hom}_{k-2}$ and the kernel has
dimension exactly $(k+1)^2$:
\[
\begin{array}{c|ccccccccc}
k & 0 & 1 & 2 & 3 & 4 & 5 & 6 & 7 & 8\\ \hline
\lambda_k = k(k+2) & 0 & 3 & 8 & 15 & 24 & 35 & 48 & 63 & 80\\
\text{multiplicity } (k+1)^2 & 1 & 4 & 9 & 16 & 25 & 36 & 49 & 64 & 81
\end{array}
\]
Every nonzero eigenvalue has multiplicity $> 2$; already $\lambda_1$ has
multiplicity $4$.
\end{proof}

\section{The conjecture's own example family refutes its own bound}

The conjecture proposes ``the family of rotated round metrics'' as the
counterexample to multiplicity rigidity. But a rotation
$Q \in SO(4)$ is an isometry of $(\SSp^3, g_r)$: it preserves round distances
($Q^T Q = I$), so the pulled-back metric $Q^* g_r$ is \emph{identically}
$g_r$. The ``family of rotated round metrics'' is a single isometry class
carrying a single spectrum, with $\mathrm{mult}(\lambda_k) = (k+1)^2$ on
every member --- in particular $\mathrm{mult}(\lambda_1) = 4 > 2$ on every
member. So the proposed family does not merely fail to support the claimed
upper bound of two; it is precisely a set of witnesses against it. The two
clauses of the conjecture are jointly inconsistent:

\begin{itemize}
\item the bound clause (``multiplicity upper bound is two'') is refuted by
$g_r$ and by every rotated round metric, all of which are Zoll metrics on
$\mathrm{SU}(2)$;
\item the counterexample clause is vacuous as stated, because the proposed
counterexample family is one isometry class with no metric-to-metric
spectrum variation to exhibit, and it contradicts the bound clause.
\end{itemize}

A single counterexample suffices to refute the universal bound: the round
metric $g_r$ is a Zoll metric on $\mathrm{SU}(2)$ with
$\mathrm{mult}(\lambda_1) = 4$.

\section{Boundary remarks}

\begin{itemize}
\item \emph{Scale invariance.} The round metric of any radius $r$ has
eigenvalues $k(k+2)/r^2$ with the same multiplicities $(k+1)^2$, so no
rescaling escapes the counterexample. More generally every bi-invariant
metric on $\mathrm{SU}(2)$ is isometric to a round $\SSp^3$.
\item \emph{Formalization.} The discrete skeleton is machine-checked in pure
Lean~4 (no Mathlib, no axioms, no \texttt{sorry}; see \texttt{lean4/}): the
eigenvalue and multiplicity tables, the dimension count
$\#\mathrm{Hom}_k - \#\mathrm{Hom}_{k-2} = (k+1)^2$ for $k \le 12$, the
explicit $k = 1$ case, the linear independence of the coordinate functions
via the identity evaluation matrix, and the rotation isometry
(norm preservation and permutation of the axis points).
\item \emph{Generically Zoll metrics may well have simple spectrum} (as
expected on $S^3$ by analogy with negative-curvature results); that is
irrelevant to the refuted statement, which is a \emph{universal upper bound}
claim falsified by an explicit Zoll metric the conjecture itself points at.
\end{itemize}

\end{document}
